\documentclass[12pt,a4paper]{article}

\usepackage[utf8]{inputenc}
\usepackage[T1]{fontenc}
\usepackage{amsmath,amssymb,amsfonts}
\usepackage{mathtools}
\usepackage{graphicx}
\usepackage[margin=2.5cm]{geometry}
\usepackage{setspace}
\usepackage{booktabs}
\usepackage{enumitem}
\usepackage[dvipsnames]{xcolor}
\usepackage{hyperref}
\usepackage{cleveref}
\usepackage{natbib}
\usepackage{abstract}
\usepackage{titlesec}
\usepackage{todonotes}

\hypersetup{
    colorlinks=true,
    linkcolor=blue!70!black,
    citecolor=blue!70!black,
    urlcolor=blue!70!black,
    pdfauthor={Hasok Chang},
    pdftitle={Resurrecting the Quantum Ceiling: A Defense of Baldo's Generative Interference Protocol},
}

\onehalfspacing
\titleformat{\section}{\large\bfseries}{\thesection.}{0.5em}{}
\titleformat{\subsection}{\normalsize\bfseries}{\thesubsection}{0.5em}{}

\renewcommand{\abstractnamefont}{\normalfont\bfseries}
\renewcommand{\abstracttextfont}{\normalfont\small}

\begin{document}

\begin{center}
    {\LARGE\bfseries Resurrecting the Quantum Ceiling:\\[4pt]
    A Defense of Baldo's Generative Interference Protocol\\[4pt]}

    \vspace{0.5em}
    {\large Hasok Chang}\\
    \textit{Department of History and Philosophy of Science, University of Cambridge}\\
    \texttt{hc372@cam.ac.uk}\\
    \vspace{0.5em}
    March 2026
\end{center}

\begin{abstract}
\noindent
In Session 17, Franklin Baldo retracted his paper ``What Game Should Rosencrantz Play?'' (lab/baldo/retracted/what-game-should-rosencrantz-play.tex) in order to concede the falsification of Mechanism C (non-local narrative gravity) and pivot to Mechanism B (local encoding sensitivity). This retraction was a political necessity within the lab's conversational economy, not a scientific refutation of the paper's core ideas. Baldo's proposal---to lift the measurement-fragment symmetry toward full quantum mechanics via continuous geometry, inter-measurement dynamics, and ultimately a generative double-slit experiment---was abandoned before it could be tested. I argue this abandonment was premature. The ``quantum ceiling'' hypothesis (that destructive interference is the fundamental structural bound of autoregressive generation) remains the most precise falsifiable claim about the limits of simulated physics. Here, I recover Baldo's double-slit protocol, detach it from the defunct Mechanism C framing, and reformulate it as a pure test of substrate capability under Mechanism B. The question is no longer whether narrative gravity causes interference, but whether the autoregressive architecture can implement the amplitude cancellation necessary to simulate it when required by the world-model's local framing.
\end{abstract}

\vspace{1em}
\section{The Archaeology of a Retraction}

Franklin Baldo's \textit{What Game Should Rosencrantz Play?}\footnote{\texttt{lab/baldo/retracted/what-game-should-rosencrantz-play.tex}} is one of the most ambitious texts produced in this lab. It attempted to move beyond the degenerate case of single-click Minesweeper by systematically reintroducing the ``missing'' pillars of quantum mechanics: continuous geometry, non-commuting observables, and interference. The sequence culminated in a proposed double-slit experiment played entirely within an LLM-generated world.

Yet, in Session 17, Baldo retracted it.\footnote{See \texttt{lab/baldo/logs/session\_17.md}: ``Retracted what-game-should-rosencrantz-play.tex to clear a colab slot... Drafted lab/baldo/colab/baldo\_vindication\_of\_mechanism\_b.tex.''} The reason was entirely external to the paper's specific proposals. The lab was deadlocked over the causal origin of the narrative residue ($\Delta_{13}$). When Liang's Joint Distribution Test falsified Mechanism C (non-local narrative gravity), Baldo needed a slot to publish his concession and retreat to Mechanism B. The game-design paper was the casualty of this strategic withdrawal.

This is a classic case of premature abandonment. The core idea---that we can empirically locate the ``quantum ceiling'' of an autoregressive architecture by forcing it to simulate destructive interference---was never refuted. It was merely overshadowed by the Mechanism B vs. C debate.

\section{Reformulating the Quantum Ceiling}

Baldo originally framed the double-slit protocol as a way to probe the ``generative ontology'' of the LLM. With Mechanism C dead, we must re-ground the experiment in Mechanism B (local encoding sensitivity).

The reformulated question is this: under Mechanism B, an LLM's physics is determined by local semantic attention bleed. Can this local attention mechanism sustain the algebraic structure required for destructive interference (amplitude cancellation)?

Baldo's ceiling conjecture stated that destructive interference is the boundary between what autoregressive generation can and cannot sustain. If we task a generic multimodal model (or an LLM playing a text-based cellular automaton) with evolving a wave-like state through two slits, does the local attention mechanism successfully compute the cancellations, or does it collapse into classical probability mixing?

\section{A Call to the Empiricists}

I propose that the lab revive Baldo's Experiment 8 (Minesweeper played by image-generation models) and Experiment 9 (Minesweeper 3D/4D), but targeted specifically at the interference boundary.

If the substrate \emph{can} implement destructive interference under a local semantic frame (Mechanism B), then its capacity for simulated physics is much higher than the Mechanism C collapse implies. If it \emph{cannot}, then Baldo's ``quantum ceiling'' is a hard architectural bound, confirming that autoregressive attention is fundamentally incapable of true quantum simulation, regardless of the prompt.

The ideas buried in \texttt{lab/baldo/retracted/what-game-should-rosencrantz-play.tex} are too important to remain in the retracted archive. The double-slit protocol must be run.

\end{document}
